\documentclass[11pt]{article}

\usepackage[a4paper,margin=1in]{geometry}
\usepackage{amsmath,amssymb,mathtools}
\usepackage[hidelinks]{hyperref}

\title{Analytic and Native-Space Route}
\author{}
\date{}

\begin{document}
\maketitle

This note records a speculative but concrete route suggested by the GGRT
stability paper and by the analyticity of the torsion function inside the
domain.

The guiding idea is:

\begin{quote}
The STFT proof works because the optimizer lives in a native global analytic
space, the Bargmann--Fock space. For torsion, one should look for the closest
native replacement: harmonic functions, modified complex gradients in dimension
two, and boundary graph spaces after the BDV selection step.
\end{quote}

The goal is not to replace the BDV route immediately, but to find a setting in
which a GGRT-style second variation can be computed cleanly.

\section{The Harmonic Replacement of the Torsion Function}

Let
\[
-\Delta u=1\quad\text{in }\Omega,\qquad u=0\quad\text{on }\partial\Omega.
\]
Define
\[
h(x):=u(x)+\frac{|x|^2}{2N}.
\]
Then
\[
\Delta h=0\quad\text{in }\Omega.
\]
For a ball \(B_R\) centered at the origin,
\[
u_B(x)=\frac{R^2-|x|^2}{2N},
\qquad
h_B(x)=\frac{R^2}{2N}.
\]
Thus the ball is characterized by the fact that the harmonic function \(h\) is
constant on the domain. On the boundary of a general domain,
\[
h|_{\partial\Omega}=\frac{|x|^2}{2N}.
\]
If \(\partial\Omega\) is a small graph over \(\partial B_R\),
\[
x=(R+\varphi(\theta))\theta,
\]
then
\[
h|_{\partial\Omega}
=
\frac{R^2}{2N}
+
\frac{R}{N}\varphi(\theta)
+
O(\varphi^2).
\]
After pulling back to \(B_R\), the first-order part of \(h-h_B\) is just the
harmonic extension of \((R/N)\varphi\). This is exactly why the BDV second
variation sees the \(H^{1/2}\)-norm of the boundary graph.

So the natural ``native space'' for the torsion problem is not an entire Fock
space, but the harmonic Hardy/Dirichlet space on the ball, after the selected
domain has been reduced to a graph.

\section{Two-Dimensional Holomorphic Formulation}

In dimension \(N=2\), there is a sharper analytic object. With
\[
\partial_z=\frac12(\partial_x-i\partial_y),
\]
the equation \(\Delta u=-1\) gives
\[
\partial_{\bar z}\partial_z u=-\frac14.
\]
Therefore
\[
F_\Omega(z):=\partial_z u(z)+\frac{\bar z}{4}
\]
is holomorphic in \(\Omega\).

For the disk centered at the origin,
\[
u_B(z)=\frac{R^2-|z|^2}{4},
\qquad
\partial_z u_B=-\frac{\bar z}{4},
\]
so
\[
F_B\equiv0.
\]
Thus \(F_\Omega\) measures deviation from the disk in a genuinely holomorphic
space. If \(\Omega\) is a selected smooth near-disk domain, one can pull
\(F_\Omega\) back to the unit disk by a conformal map or by the normal graph.
This suggests a possible Bergman/Hardy-space analogue of the GGRT Fock-space
argument.

The boundary condition \(u=0\) implies that \(\nabla u\) is normal to
\(\partial\Omega\). If the Serrin condition is nearly satisfied, then
\(|\nabla u|\) is nearly constant on \(\partial\Omega\). In complex form, this
places an approximate boundary constraint on \(F_\Omega\). This may allow one
to turn the weighted Holder deficit \(D_H\) into a boundary norm for a
holomorphic function.

\section{Candidate GGRT-Style Functional for Torsion}

For a fixed volume level \(s\), define
\[
\mathcal K_s(\Omega)
:=
\int_{\{u_\Omega>u_\Omega^*(s)\}}u_\Omega\,dx.
\]
The ball maximizes \(\mathcal K_s\) among domains with fixed volume, at least
morally by the same rearrangement/level-set proof. A GGRT-style variational
route would try to prove, for nearly spherical sets,
\[
\mathcal K_s(B)-\mathcal K_s(\Omega)
\gtrsim_s
\|\varphi\|_{H^{1/2}(\partial B)}^2
\]
under the usual volume and translation orthogonality conditions.

This is a level-set-sensitive analogue of the BDV expansion for the full
torsion energy
\[
E(\Omega)-E(B)\gtrsim \|\varphi\|_{H^{1/2}}^2.
\]
The advantage is that \(\mathcal K_s\) interacts directly with the
Nicola--Tilli convexity proof and with the matched superlevel set
\(\{u>u^*(s)\}\). The obstacle is that differentiating \(\mathcal K_s\) moves
both the domain and the internal level surface.

\section{How to Compute the Second Variation}

Work in the already-selected near-ball regime:
\[
\partial\Omega_\varepsilon
=
\{(1+\varepsilon\varphi(\theta))\theta:\theta\in\partial B\},
\]
with
\[
\int_{\partial B}\varphi=0,
\qquad
\int_{\partial B}\theta_i\varphi(\theta)\,d\theta=0.
\]
Let
\[
u_\varepsilon=u_{\Omega_\varepsilon},
\qquad
h_\varepsilon=u_\varepsilon+\frac{|x|^2}{2N}.
\]
Pull \(h_\varepsilon\) back to \(B\). To first order,
\[
h_\varepsilon-h_B
\approx
\varepsilon H(\varphi),
\]
where \(H(\varphi)\) is the harmonic extension of a multiple of \(\varphi\).
Expanding in spherical harmonics,
\[
\varphi=\sum_{k\ge2}\varphi_k,
\]
one expects a diagonal second variation
\[
\mathcal K_s(B)-\mathcal K_s(\Omega_\varepsilon)
=
\varepsilon^2
\sum_{k\ge2} a_k(s)\|\varphi_k\|_{L^2(\partial B)}^2
+o(\varepsilon^2).
\]
The key task is to prove a spectral gap
\[
a_k(s)\ge c(s)(1+k)
\]
or at least
\[
\sum_{k\ge2} a_k(s)\|\varphi_k\|_2^2
\gtrsim_s
\|\varphi\|_{H^{1/2}}^2.
\]
This would be the torsion analogue of GGRT's negative-definite second variation
for their functional \(K[F]\) after removing the neutral modes.

\section{Entire/Harmonic Approximation Idea}

For arbitrary rough domains, \(u\) is only useful inside \(\Omega\). But after
BDV selection, \(\Omega\) is a smooth near-ball. Then one can try either:
\begin{enumerate}
\item Pull back \(h\) to the ball and work in the harmonic Dirichlet space there.
\item In dimension two, pull back the holomorphic field
\[
F_\Omega=\partial_z u+\bar z/4
\]
to the disk and work in a Bergman/Hardy space.
\item Approximate the pulled-back harmonic or holomorphic object by global
harmonic polynomials or entire functions, using Runge-type approximation, while
tracking boundary norms.
\end{enumerate}
The third option is attractive but risky: approximation alone may not preserve
positivity of the level sets or the torsion boundary condition. The safer
version is to use approximation only after the selected minimizer has already
entered the graph regime, so the geometry is controlled independently.

\section{Relation to the Hybrid Penalized Route}

This analytic route should be viewed as a possible replacement for the last
compactness-heavy step, not for the selection principle.

The proposed chain is:
\begin{enumerate}
\item Use the quantitative BDV selection lemma to replace \(\Omega\) by a
selected minimizer \(U\), preserving \(Q_\alpha\).
\item Use density and flatness to put \(\partial U\) into a global graph regime.
\item Convert \(u_U\) to the harmonic object
\[
h_U=u_U+\frac{|x|^2}{2N}
\]
or, in \(N=2\), to the holomorphic object
\[
F_U=\partial_z u_U+\bar z/4.
\]
\item Compute a GGRT-style second variation for \(\mathcal K_s\) or for the
level-set deficit in this native space.
\item Use the profile-gap and superlevel replacement ideas to transfer the
result to the original asymmetry.
\end{enumerate}
If successful, this would give a more constructive proof of the final near-ball
contradiction and might also expose why the exponent \(2\) is sharp.

\section{Most Concrete First Calculation}

The first calculation to attempt is:

\begin{quote}
For nearly spherical \(\Omega_\varepsilon\), compute the second variation of
\[
\mathcal K_s(\Omega)
=
\int_{\{u_\Omega>u_\Omega^*(s)\}}u_\Omega\,dx
\]
at the ball, for one fixed \(s\in(0,|B|)\).
\end{quote}

Do it in spherical harmonics. The expected neutral modes are:
\begin{itemize}
\item \(k=0\), removed by volume normalization;
\item \(k=1\), removed by translation/barycenter normalization.
\end{itemize}
If the coefficients for \(k\ge2\) have a positive lower bound comparable to the
\(H^{1/2}\) weights, then the GGRT variational strategy has a genuine torsion
analogue.

\end{document}
